\section{The Ventriloquist and the Dummy}

Despite appearances, this article is the prequel to a post in the Sanatana Dharma.

At a small prestigious university just outside of Boston, I was first exposed to the ideas of the New Left. The goal of Vatican II was to make the faith intelligible to modern man, but its effect was to make it irrelevant. If you think like a modern, then why not embrace it fully?

Left to my own devices, I attempted to meld a worldview out of elements from \textbf{Alan Watts}, \textbf{Timothy Leary}, \textbf{Carlos Castaneda} on the one hand and the New Left on the other. I was to explore those ideologies in some depth. The benefit of that exercise is that I understand the assumptions and forms of logic associated with those movements, so I am sensitive to their effects. That is unlike most of my cohort who have not had a new idea since the age of 19, although they are materially intelligent and nominally educated. Hence, they are dead-certain about their opinions with no real understanding of their intellectual sources.

I was a reader of \emph{Ramparts} magazine and the \emph{Guardian} newspaper. Although turgid, the analyses presented in the latter were detailed and logically coherent, at least as I now remember them. That is still a model to follow.

Of course, the master at that time was \textbf{Herbert Marcuse}, the ``Evola of the left'' (although Evola was utterly invisible at that time). It would be a mistake to underestimate his intelligence and prescience, and I still feel his influence. His accomplishment was to reconnect Marxism with German Idealism, existentialism, and more famously, with Freudian psychology. There were three lessons in particular that have been long lasting. It is always a benefit to wrestle with serious thought, even if one ultimately disagrees with it.

In \emph{Reason and Revolution}, Marcuse made the case that rationality cannot be restricted to mathematics and positive science. At that time positivism was a dominant philosophical movement and I simply assumed that physics was the only rational approach to the world. Idealism presumes that ideas or consciousness transcend empirical thought. That provided the intellectual cover to embrace someone like an Alan Watts. I preferred his earlier books since they were an attempt to relate Christianity to the \emph{philosophia perennis}; it was there that I first read the name Rene Guenon, although it would be quite some time before I became acquainted with his actual works.

The next lesson was that Marcuse provided a link from the thought of the community and the individual man's consciousness. Marxism itself had no psychology, so the process from social change to the mental states of the individual was left unclear. I saw Marcuse as demonstrating a continuity from the group to the individual; that is something I still accept.

Unfortunately, Marcuse used the Freudian model for man's psychology. In a sense, it is a reworking and a distortion of traditional psychology. The id was the source of instinctual impulses, particularly the libido and willfulness. Hence, it embraces both the eros and the thymos. The ego is the empirical self, but its purpose in that system is to satisfy the urges of the id. The superego is the transcendent element in the psyche. However, it is not a true transcendence but represents the internalization of the thought of the community. Therefore, by modifying the social forms and relations, the superego is altered, thereby affecting the individual, presumably ``liberating'' him.

That leads to the third lesson, which is the connection between sexual liberation and the overturning of patriarchal, hierarchal social arrangements, as described in his masterwork, \emph{Eros and Civilization}. Without going into all the details about ``surplus repression'', the premise is that technology would alleviate man's need to work. By overthrowing the patriarchal system that represses humans through the superego, man would become ``polymorphously perverse'', that is, his entire body, not just the genital region, would become sexually aware and a source of pleasure. There would then be no restriction of sexual satisfaction.

The truth of this insight is all too obvious; the more the former order of things is overturned, the more demand there is for perverse sexuality. However, this is not at all a sign of liberation as was expected and is still claimed, but rather a tool of manipulation. Once the superego has lost whatever tenuous connection it may have had with the cosmic order, then almost any thoughts can be imprinted in the minds of those addicted to sexual expression. The reason is obvious. The overthrow of the cosmic order does not result in a new order of things, but rather in a lack of order, or chaos. This is nothing new, whatever you may have been told in certain circles. For a cultural history of the use of sex as a form of social control, I recommend the works of E Michael Jones.

\paragraph{Drinking for Free}

After graduation, I may have been able to continue along those lines had I entered a ``soft'' field like education, psychology, politics, or government service for example. Many others did and their effects have colored that past forty years.

Instead I was involved with mathematics, actuarial science, and eventually computer technology. It was a world of work unlike that described in the periodicals I used to work. The effort necessary to move people and matter is quite different from manipulating words on a page.

Beyond that, dating suddenly became much more formal, and there were no longer hippy chicks available to perpetuate the former lifestyle. You first had to call, you had to plan a dinner, etc. I recall early on meeting someone at a club. When I called her a few days later and heard her infant crying in the background, I abruptly hung up on her. Unwed motherhood was still not socially acceptable, even for a wannabe radical.

At that time the disco era was peaking along with the popularity of backgammon. Different clubs would have a backgammon tournament before the music started, with prize money running from \$50 to \$100 or so. Unlike today, you could have a good time at a club with that much cash. Since there are few things more satisfying that getting a drink for free, and I was pretty skilled at the game, I was sure to show up for those contests. More often than not, I would get some of that cash, and it was like drinking for free.

There was one tournament that I recall being paired against an older woman from New York. Falling behind, I made the only move that would give me any chance for victory. On my next roll I hit the double-sixes and took the game. She became visibly annoyed and pointed out that I had made a really bad move. Unfortunately for her case, unlike a sport like gymnastics, you don't get any ``style'' points in backgammon.

She had no idea how to deal with the goddess Fortuna. She probably read a book on the game or took a lesson at her condo clubhouse. In vain did she expect that by following some rules or guidelines she could control Fortuna. Not at all, although you must never put the Lord, thy God, to the test, you should always test Fortuna.

\paragraph{The Last Dummy}


Besides backgammon, there were other ways to win money at a club. One club had a talent night with a cash prize, so I decided on a gambit. I approached an attractive young woman and convinced her to enter the contest as a faux ventriloquist act. I would play the ventriloquist and she the dummy. She would just have to sit on my lap while I stuck my hand up her back and pretended to pull a string behind her head. She would just have to open and close her mouth and move her head around as I told jokes.

She went for the idea except she insisted that she be the ventriloquist and I the dummy. I did what I could, but she wasn't funny, so we got yanked off the stage. When I phoned her a few days later a man answered and called her to the phone. In the background I could hear voices and laughter, something about the ``guy from the talent show''. She was non-committal about a date and seemed reticent to talk. I sensed I was being taken for the dummy. Since there was no move that could result in victory no matter how the dice rolled, I resolved that my life would have to take a more serious turn.


\flrightit{Posted on 2013-10-28 by Cologero}

\begin{center}* * *\end{center}

\begin{footnotesize}
\begin{sffamily}
\texttt{Avery on 2013-10-28 at 12:06 said:}

Confucius said, ``Someone who spends all his time fulfilling his appetites and never once exercises his mind is a hard case indeed. There are games of luck and chance. It would be better to play those than to do nothing.'' (17.22)

For those unfamiliar with Herbert Marcuse, I wrote about him here:

\url{http://avery.morrow.name/blog/2012/07/herbert-marcuse-the-original-tumblard/}


\end{sffamily}
\end{footnotesize}

